\section{Solución propuesta}

La solución propuesta es \textbf{SocioUnido}, una plataforma de \textit{software como servicio (Software as a Service - SaaS)} con enfoque de empresa a empresa (B2B) que centraliza la gestión administrativa, financiera y social de los clubes de fútbol. El sistema actúa como un ecosistema integral que permite a las instituciones optimizar su recaudación, fidelizar a su masa societaria de manera proactiva y segurizar los accesos a los estadios mediante tecnología de vanguardia.

La plataforma se estructura en torno a los siguientes componentes clave:

\begin{itemize}
    \item \textbf{Panel administrativo (Web):} panel de control integral de escritorio diseñado para la dirigencia y tesorería. Permite visualizar métricas complejas, consolidar la recaudación y gestionar campañas institucionales.
    \item \textbf{Aplicación (App) móvil (Aplicación web progresiva [PWA]):}
    \begin{itemize}
        \item \textbf{Interfaz de socio:} interfaz liviana y accesible para el hincha. Provee una billetera digital que aloja el carnet dinámico, facilita el pago de cuotas y permite la reserva de espacios o servicios del club.
        \item \textbf{Interfaz de personal:} aplicación móvil de uso exclusivo para el personal de seguridad y control en los accesos, optimizada para el escaneo de códigos QR a alta velocidad.
    \end{itemize}
    \item \textbf{Acceso Inteligente (Control sin conexión):} módulo de seguridad perimetral basado en la generación de vales (tokens) de contraseñas de un solo uso basadas en el tiempo (TOTP por sus siglas en inglés, \textit{Time-Based One-Time Password}). Permite validar matemáticamente los accesos en escenarios de alta concurrencia sin depender de la conectividad a internet.
    \item \textbf{Asistente conversacional (Procesamiento de Lenguaje Natural [PLN]):} integración omnicanal vía WhatsApp que procesa el lenguaje natural del usuario para resolver consultas administrativas y orquestar transacciones de pago automáticas.
    \item \textbf{Motor predictivo (IA y analítica [analytics]):} subsistema estadístico que analiza el historial de pagos y asistencia para calcular el riesgo de morosidad (abandono) de cada socio, disparando alertas y campañas de retención tempranas.
\end{itemize}

\subsection*{Tecnologías a evaluar}

Si bien el conjunto de tecnologías (\textit{stack}) tecnológico definitivo se encuentra actualmente en proceso de definición, las principales opciones que se están evaluando para soportar la arquitectura orientada a microservicios son:

\begin{itemize}
    \item \textbf{Parte frontal (Frontend) y aplicaciones móviles:} React, HTML5 y CSS3 para el desarrollo del panel de control (\textit{Dashboard}) web y la aplicación web progresiva (\textit{Progressive Web App} - PWA) del socio. React Native para el desarrollo nativo de la Aplicación (App) de Validación, priorizando el rendimiento en el escaneo óptico.
    \item \textbf{Parte trasera [Backend] (Microservicios):} lenguaje Python estructurado sobre el marco de trabajo (\textit{framework}) FastAPI para los módulos de Gestión (Principal o Core), Autenticación, Pagos, Procesamiento de Lenguaje Natural (PLN) y Analítica (Analytics). Go (Golang) para el microservicio crítico de acceso inteligente (\textit{Smart Access}), garantizando máxima velocidad de validación y manejo de concurrencia.
    \item \textbf{Puerta de enlace de API (API Gateway):} adopción de KrakenD como puerta de enlace para recibir, validar mediante vales web JSON (JWT) y enrutar todas las peticiones hacia los microservicios correspondientes.
    \item \textbf{Bases de datos:} PostgreSQL como motor relacional principal para asegurar la integridad transaccional del padrón de socios y los historiales financieros. Redis como almacenamiento en memoria (caché) para el módulo contra ataques de repetición (anti-replay) en los accesos. Firebase para la gestión de credenciales y roles.
    \newpage
    \item \textbf{Inteligencia artificial:} Scikit-Learn y Pandas para el modelado estadístico y tareas asíncronas de predicción de abandono (\textit{ChurnPredictionJob}); LangChain como motor de procesamiento de lenguaje natural (PLN) para el reconocimiento de intenciones en el canal conversacional.
    \item \textbf{Integraciones de terceros:} implementación de kits de desarrollo de software (SDKs) para pasarelas de pago externas (ej. MercadoPago) y consumo de retrollamadas web (\textit{webhooks}) públicas mediante la interfaz de programación de aplicaciones (API) oficial de WhatsApp (Meta).
\end{itemize}

Estas decisiones preliminares se encuentran estrictamente guiadas por criterios de escalabilidad, facilidad de integración modular, tolerancia a fallos en estadios y soporte a largo plazo. Cabe destacar que el diseño arquitectónico detallado planteado en esta instancia de ideación se encuentra plasmado gráficamente en el \textbf{Anexo A: Modelo de arquitectura C4} del presente documento.
